% Created 2025-01-20 lun 14:25
% Intended LaTeX compiler: pdflatex
\documentclass[11pt]{article}
\usepackage[utf8]{inputenc}
\usepackage[T1]{fontenc}
\usepackage{graphicx}
\usepackage{longtable}
\usepackage{wrapfig}
\usepackage{rotating}
\usepackage[normalem]{ulem}
\usepackage{amsmath}
\usepackage{amssymb}
\usepackage{capt-of}
\usepackage{hyperref}
\usepackage{minted}

\usepackage{parskip}
\usepackage[utf8]{inputenc} %% For unicode chars
\usepackage{placeins}
\usepackage[margin=2.50cm]{geometry}
\usepackage[T1]{fontenc}
\usepackage{mathpazo}
\linespread{1.05}
\usepackage[scaled]{helvet}
\usepackage{courier}
\hypersetup{colorlinks=true,linkcolor=blue}
\RequirePackage{fancyvrb}
\DefineVerbatimEnvironment{verbatim}{Verbatim}{fontsize=\small,formatcom = {\color[rgb]{0.5,0,0}}}
\usepackage{mdframed}
\BeforeBeginEnvironment{minted}{\begin{mdframed}}
\AfterEndEnvironment{minted}{\end{mdframed}}
\setcounter{secnumdepth}{3}
\author{José Manuel Sánchez Álvarez}
\date{}
\title{PHP and Databases}
\hypersetup{
 pdfauthor={José Manuel Sánchez Álvarez},
 pdftitle={PHP and Databases},
 pdfkeywords={},
 pdfsubject={},
 pdfcreator={Emacs 27.1 (Org mode 9.6.18)}, 
 pdflang={Spanish}}
\begin{document}

\maketitle
\setcounter{tocdepth}{3}
\tableofcontents

\setlength\parindent{10pt}

\section{PHP and Databases}
\label{sec:org263cae5}
Working with databases is a critical aspect of PHP development. In this
section, we'll explore how PHP interacts with databases, particularly
focusing on MySQL and PostgreSQL.

\subsection{Introduction to Databases}
\label{sec:org5b7275e}
A database is an organized collection of data, generally stored and
accessed electronically from a computer system. In web development,
databases are used to store data such as user information, posts,
comments, and more.

\subsection{Database Access from PHP}
\label{sec:org5ab1b16}
PHP interacts with databases through specific libraries. The most common
libraries for MySQL are:

\begin{itemize}
\item \textbf{MySQLi}: An improved version of the older MySQL library. It provides
both procedural and object-oriented interfaces and supports prepared
statements, which are crucial for security.

\item \textbf{PDO (PHP Data Objects)}: A database access layer providing a uniform
method of access to multiple databases. It doesn't support
database-specific syntax, but it's flexible and supports prepared
statements.
\end{itemize}

\subsubsection{MySQLi and PDO for MySQL}
\label{sec:org78cc2f3}
\begin{itemize}
\item Connecting to a MySQL Database
\label{sec:org22497bb}
\begin{itemize}
\item \textbf{MySQLi}:

\begin{minted}[]{php}
$mysqli = new mysqli("localhost", "username", "password", "database");

if ($mysqli->connect_error) {
    die("Connection failed: " . $mysqli->connect_error);
}
\end{minted}

\item \textbf{PDO}:

\begin{minted}[]{php}
try {
    $pdo = new PDO("mysql:host=localhost;dbname=database",
                                      "username", "password");
    // Set the PDO error mode to exception
    $pdo->setAttribute(PDO::ATTR_ERRMODE, PDO::ERRMODE_EXCEPTION);
} catch(PDOException $e) {
    die("Connection failed: " . $e->getMessage());
}
\end{minted}
\end{itemize}

\item Listing Databases
\label{sec:org1a90090}
\begin{itemize}
\item \textbf{MySQLi}:

\begin{minted}[]{php}
$result = $mysqli->query("SHOW DATABASES");

while ($row = $result->fetch_assoc()) {
    echo $row["Database"] . "\n";
}
\end{minted}

\item \textbf{PDO}:

\begin{minted}[]{php}
$query = $pdo->query("SHOW DATABASES");

while ($row = $query->fetch(PDO::FETCH_ASSOC)) {
    echo $row["Database"] . "\n";
}
\end{minted}
\end{itemize}
\end{itemize}

\subsubsection{CRUD Operations}
\label{sec:org9e337bf}
\begin{itemize}
\item Query All Data from a Table
\label{sec:org42ac7f1}
\begin{itemize}
\item \textbf{MySQLi}:

\begin{minted}[]{php}
$result = $mysqli->query("SELECT * FROM table_name");

while ($row = $result->fetch_assoc()) {
    echo $row["column1"] . ", " . $row["column2"] . "\n";
}
\end{minted}

\item \textbf{PDO}:

\begin{minted}[]{php}
$stmt = $pdo->query("SELECT * FROM table_name");

while ($row = $stmt->fetch(PDO::FETCH_ASSOC)) {
    echo $row["column1"] . ", " . $row["column2"] . "\n";
}
\end{minted}
\end{itemize}

\item More Complex Queries
\label{sec:orgb411ebf}
You can perform complex SQL queries such as JOINs, nested SELECTs, etc.,
using the same methods.

\item Insert Data
\label{sec:org2b353e3}
\begin{itemize}
\item \textbf{MySQLi}:

\begin{minted}[]{php}
$sql = "INSERT INTO table_name (column1, column2)
  VALUES ('value1', 'value2')";
$mysqli->query($sql);
\end{minted}

\item \textbf{PDO}:

\begin{minted}[]{php}
$sql = "INSERT INTO table_name (column1, column2)
  VALUES ('value1', 'value2')";
$pdo->exec($sql);
\end{minted}
\end{itemize}

\item Delete Rows
\label{sec:orgaefc8f3}
\begin{itemize}
\item \textbf{MySQLi}:

\begin{minted}[]{php}
$sql = "DELETE FROM table_name WHERE condition";
$mysqli->query($sql);
\end{minted}

\item \textbf{PDO}:

\begin{minted}[]{php}
$sql = "DELETE FROM table_name WHERE condition";
$pdo->exec($sql);
\end{minted}
\end{itemize}

\item Create Tables
\label{sec:orgee2bf01}
\begin{itemize}
\item \textbf{MySQLi}:

\begin{minted}[]{php}
$sql = "CREATE TABLE table_name (id INT, data VARCHAR(100))";
$mysqli->query($sql);
\end{minted}

\item \textbf{PDO}:

\begin{minted}[]{php}
$sql = "CREATE TABLE table_name (id INT, data VARCHAR(100))";
$pdo->exec($sql);
\end{minted}
\end{itemize}
\end{itemize}

\subsubsection{PostgreSQL Connection and CRUD Operations with PDO}
\label{sec:org709bbec}
Connecting to PostgreSQL and performing CRUD operations with PDO is
similar to MySQL. The primary difference lies in the DSN (Data Source
Name) used in the PDO constructor.

\begin{itemize}
\item Connecting to a PostgreSQL Database
\label{sec:orgbafb765}
\begin{minted}[]{php}
try {
    $pdo = new PDO("pgsql:host=localhost;dbname=database",
                                     l"username", "password");
    $pdo->setAttribute(PDO::ATTR_ERRMODE, PDO::ERRMODE_EXCEPTION);
} catch(PDOException $e) {
    die("Connection failed: " . $e->getMessage());
}
\end{minted}

The CRUD operations for PostgreSQL using PDO are identical to those for
MySQL, as PDO abstracts the database details, providing a consistent
method for database interaction.
\end{itemize}

\subsection{CRUD Car Dealership}
\label{sec:org61cae48}
Creating a complete set of CRUD operations for a car dealership database
involves several components, including database design, PHP scripts for
handling operations, and HTML forms for user input. Let's outline a
basic structure for such an application.

\subsubsection{Database Structure}
\label{sec:orgabe3ddc}
\begin{enumerate}
\item \textbf{Tables:}

\begin{itemize}
\item \texttt{Car}: \texttt{id}, \texttt{model}, \texttt{price}, \texttt{brand\_id}, \texttt{seller\_id}
\item \texttt{CarBrand}: \texttt{id}, \texttt{name}
\item \texttt{Seller}: \texttt{id}, \texttt{name}, \texttt{cars\_selled}
\item \texttt{Client}: \texttt{id}, \texttt{name}
\item \texttt{Purchase}: \texttt{id}, \texttt{client\_id}, \texttt{car\_id}, \texttt{date}
\end{itemize}

\item \textbf{Relationships:}

\begin{itemize}
\item A \texttt{Car} belongs to a \texttt{CarBrand} and a \texttt{Seller}.
\item A \texttt{Purchase} links a \texttt{Client} to a \texttt{Car}.
\end{itemize}
\end{enumerate}

\subsubsection{HTML Forms for Data Insertion}
\label{sec:org8eb23f9}
\begin{enumerate}
\item \textbf{Add a New Car:}

\begin{minted}[]{html}
<form action="add_car.php" method="post">
    Model: <input type="text" name="model"><br>
    Price: <input type="text" name="price"><br>
    Brand: <select name="brand_id">
        <!-- Options populated from CarBrand table -->
    </select><br>
    Seller: <select name="seller_id">
        <!-- Options populated from Seller table -->
    </select><br>
    <input type="submit" value="Add Car">
</form>
\end{minted}

\item \textbf{Add a New Client:}

\begin{minted}[]{html}
<form action="add_client.php" method="post">
    Name: <input type="text" name="name"><br>
    <input type="submit" value="Add Client">
</form>
\end{minted}
\end{enumerate}

\subsubsection{PHP Scripts for CRUD Operations}
\label{sec:orgacb2d51}
\begin{enumerate}
\item \textbf{Add a New Car (\texttt{add\_car.php}):}

\begin{minted}[]{php}
<?=
// Assume PDO $pdo is already connected to the database
if ($_SERVER["REQUEST_METHOD"] == "POST") {
    $stmt = $pdo->prepare(
      "INSERT INTO Car (model, price, brand_id, seller_id)
      VALUES (:model, :price, :brand_id, :seller_id)");
    $stmt->execute([
        ':model' => $_POST['model'],
        ':price' => $_POST['price'],
        ':brand_id' => $_POST['brand_id'],
        ':seller_id' => $_POST['seller_id']
    ]);
    echo "Car added successfully.";
}
\end{minted}

\item \textbf{Add a New Client (\texttt{add\_client.php}):}

\begin{minted}[]{php}
if ($_SERVER["REQUEST_METHOD"] == "POST") {
    $stmt = $pdo->prepare("INSERT INTO Client (name) VALUES (:name)");
    $stmt->execute([':name' => $_POST['name']]);
    echo "Client added successfully.";
}
\end{minted}
\end{enumerate}

\subsubsection{Listing Data}
\label{sec:org27c4347}
\begin{enumerate}
\item \textbf{List All Cars (\texttt{list\_cars.php}):}

\begin{minted}[]{php}
$stmt = $pdo->query("SELECT Car.model, Car.price, CarBrand.name AS brand, Seller.name AS seller FROM Car JOIN CarBrand ON Car.brand_id = CarBrand.id JOIN Seller ON Car.seller_id = Seller.id");
while ($row = $stmt->fetch(PDO::FETCH_ASSOC)) {
    echo "Model: " . $row['model'] . ", Price: " . $row['price'] . ", Brand: " . $row['brand'] . ", Seller: " . $row['seller'] . "<br>";
}
\end{minted}

\item \textbf{List Cars Bought by Clients (\texttt{cars\_bought\_by\_clients.php}):}

\begin{minted}[]{php}
$stmt = $pdo->query("SELECT Client.name AS client, Car.model FROM Purchase JOIN Client ON Purchase.client_id = Client.id JOIN Car ON Purchase.car_id = Car.id");
while ($row = $stmt->fetch(PDO::FETCH_ASSOC)) {
    echo "Client: " . $row['client'] . ", Car Model: " . $row['model'] . "<br>";
}
\end{minted}

\item \textbf{List Cars Sold by Sellers (\texttt{cars\_selled\_by\_sellers.php}):}

\begin{minted}[]{php}
$stmt = $pdo->query("SELECT Seller.name AS seller, COUNT(Car.id) AS cars_selled FROM Car JOIN Seller ON Car.seller_id = Seller.id GROUP BY Seller.id");
while ($row = $stmt->fetch(PDO::FETCH_ASSOC)) {
    echo "Seller: " . $row['seller'] . ", Cars Selled: " . $row['cars_selled'] . "<br>";
}
\end{minted}
\end{enumerate}

\subsubsection{Conclusion}
\label{sec:orgfcceb37}
This setup provides a basic structure for a car dealership application
with CRUD operations. The PHP scripts and HTML forms can be expanded and
customized as needed. Security considerations, like input validation and
prepared statements, have been included to prevent SQL injection. For a
real-world application, you'd also need to handle

This overview provides a foundation for working with MySQL and
PostgreSQL databases in

PHP. Each of these operations is fundamental for creating dynamic,
data-driven web applications. Understanding how to connect to a
database, perform CRUD operations, and manage different database engines
is crucial for any PHP developer.

\subsection{Prepared Statements}
\label{sec:org769d5f2}
Using prepared statements for querying databases
is an essential practice in PHP to ensure security, particularly to
prevent SQL injection attacks. Prepared statements separate SQL commands
from the data, allowing the database to recognize and treat data as
such, rather than as part of an SQL command.

\subsubsection{MySQLi and PDO}
\label{sec:orga8a7472}
PHP offers two main methods to interact with databases using prepared
statements: MySQLi (MySQL Improved) and PDO (PHP Data Objects). Both
support prepared statements, but PDO is more flexible as it can
interface with various database systems.

\begin{itemize}
\item MySQLi Prepared Statement
\label{sec:orgd9664e1}
\textbf{Connecting to the Database}:

\begin{minted}[]{php}
$mysqli = new mysqli("localhost", "username", "password", "database");
if ($mysqli->connect_error) {
    die("Connection failed: " . $mysqli->connect_error);
}
\end{minted}

\textbf{Prepared Statement}:

\begin{minted}[]{php}
// Prepare a statement
$stmt = $mysqli->prepare("SELECT * FROM table WHERE column = ?");
$stmt->bind_param("s", $param);

// Set parameters and execute
$param = "value";
$stmt->execute();

$result = $stmt->get_result();
while ($row = $result->fetch_assoc()) {
    // process $row
}

$stmt->close();
\end{minted}

\item PDO Prepared Statement
\label{sec:orgd7aefde}
\textbf{Connecting to the Database}:

\begin{minted}[]{php}
try {
    $pdo = new PDO("mysql:host=localhost;dbname=database", "username", "password");
    $pdo->setAttribute(PDO::ATTR_ERRMODE, PDO::ERRMODE_EXCEPTION);
} catch(PDOException $e) {
    die("Connection failed: " . $e->getMessage());
}
\end{minted}

\textbf{Prepared Statement}:

\begin{minted}[]{php}
$stmt = $pdo->prepare("SELECT * FROM table WHERE column = :param");
$stmt->bindParam(':param', $param, PDO::PARAM_STR);

$param = "value";
$stmt->execute();

while ($row = $stmt->fetch(PDO::FETCH_ASSOC)) {
    // process $row
}
\end{minted}
\end{itemize}

\subsubsection{CRUD Operations with Prepared Statements}
\label{sec:org12f1ebf}
\begin{itemize}
\item Insert Data
\label{sec:org19626f9}
\textbf{MySQLi}:

\begin{minted}[]{php}
$stmt = $mysqli->prepare("INSERT INTO table (column1, column2) VALUES (?, ?)");
$stmt->bind_param("ss", $column1, $column2);

$column1 = "data1";
$column2 = "data2";
$stmt->execute();
\end{minted}

\textbf{PDO}:

\begin{minted}[]{php}
$stmt = $pdo->prepare("INSERT INTO table (column1, column2) VALUES (:column1, :column2)");
$stmt->bindParam(':column1', $column1);
$stmt->bindParam(':column2', $column2);

$column1 = "data1";
$column2 = "data2";
$stmt->execute();
\end{minted}

\item Update Data
\label{sec:org81bcdf2}
\textbf{MySQLi}:

\begin{minted}[]{php}
$stmt = $mysqli->prepare("UPDATE table SET column1 = ? WHERE column2 = ?");
$stmt->bind_param("ss", $column1, $column2);

$column1 = "newdata";
$column2 = "condition";
$stmt->execute();
\end{minted}

\textbf{PDO}:

\begin{minted}[]{php}
$stmt = $pdo->prepare("UPDATE table SET column1 = :column1 WHERE column2 = :column2");
$stmt->bindParam(':column1', $column1);
$stmt->bindParam(':column2', $column2);

$column1 = "newdata";
$column2 = "condition";
$stmt->execute();
\end{minted}

\item Delete Data
\label{sec:orgee0fb4a}
\textbf{MySQLi}:

\begin{minted}[]{php}
$stmt = $mysqli->prepare("DELETE FROM table WHERE column = ?");
$stmt->bind_param("s", $param);

$param = "value";
$stmt->execute();
\end{minted}

\textbf{PDO}:

\begin{minted}[]{php}
$stmt = $pdo->prepare("DELETE FROM table WHERE column = :param");
$stmt->bindParam(':param', $param);

$param = "value";
$stmt->execute();
\end{minted}
\end{itemize}

\subsubsection{Advantages of Using Prepared Statements}
\label{sec:org94cccb6}
\begin{enumerate}
\item \textbf{Security}: They prevent SQL injection attacks.
\item \textbf{Performance}: They can improve performance, as the SQL statement
template is compiled once and can be reused with different
parameters.
\item \textbf{Flexibility}: With PDO, you can switch between different database
systems without changing the PHP code for preparing and executing the
SQL.
\end{enumerate}

\subsubsection{Conclusion}
\label{sec:org1f7e722}
Using prepared statements for database queries is a best practice in PHP
development. It enhances the security of your applications and can
improve performance. Whether using MySQLi or PDO, the concept remains
the same: separate SQL structure from the data.
\end{document}
